% ===========================================================================
%  Corps du dossier — 8 competences
% ===========================================================================

\section{Le logiciel supervise}

\begin{minipage}[t]{0.60\linewidth}
\vspace{0pt}
Le logiciel n'est ni un site vitrine ni une API transactionnelle : c'est un \textbf{jeu de role
narratif mono-joueur}, expose en REST + SSE, dont la narration est deleguee a un \textbf{fournisseur
LLM externe} et dont l'etat de partie vit \textbf{en memoire} du processus avant serialisation sur
disque. Quatre caracteristiques commandent tout le dispositif decrit ensuite.

\vspace{3pt}
\footnotesize
\begin{tabularx}{\linewidth}{@{}L{3.5cm} X@{}}
\toprule
\textbf{Caracteristique} & \textbf{Consequence} \\
\midrule
Degradation \textbf{silencieuse} prevue par conception : en cas de panne du LLM, l'application
bascule sur un narrateur local et repond toujours \code{200 OK} &
Une supervision fondee sur les codes HTTP declarerait le service sain alors que la valeur percue
s'est effondree. Une sonde du \textbf{mode de narration} est indispensable. \\
\addlinespace
Reponses en \textbf{streaming SSE} de plusieurs dizaines de secondes &
Ces routes ne peuvent pas etre jugees sur les seuils de latence des routes de jeu ; elles sont
exclues des alertes de lenteur. \\
\addlinespace
\textbf{Sessions en memoire}, sans expiration &
La consommation croit avec le nombre de joueurs depuis le demarrage : jauge de capacite. \\
\addlinespace
Fournisseur \textbf{facture a l'usage}, historique jamais tronque &
Le cout et la latence par appel croissent au fil de la partie : sonde de taille de contexte. \\
\bottomrule
\end{tabularx}
\end{minipage}\hfill
\begin{minipage}[t]{0.37\linewidth}
\vspace{0pt}
\includegraphics[width=\linewidth]{01_authentification.png}
\captionof{figure}{Authentification JWT.}
\vspace{4pt}
\includegraphics[width=\linewidth]{02_partie_en_cours.png}
\captionof{figure}{Scene d'ouverture diffusee en SSE, en mode narrateur local.}
\vspace{4pt}
\includegraphics[width=\linewidth]{07_competences.png}
\captionof{figure}{Panneau de progression : arbre de competences et points d'attribut.}
\end{minipage}

% ===========================================================================
\section{Processus de mise a jour des dependances}
\label{sec:deps}

\competence{Competence C4.1.1}{Gerer les mises a jour des dependances et des bibliotheques tiers,
en surveillant regulierement les nouvelles versions, en evaluant les impacts des mises a jour, et
en les integrant de maniere securisee.}

\subsection{Perimetre et politique}

Le perimetre couvre quatre ensembles : \textbf{backend Python} (\code{requirements.txt} :
FastAPI, Uvicorn, sse-starlette, PyJWT, bcrypt, openai, PyYAML, pytest) ; \textbf{frontend npm}
(\code{package.json} + \code{package-lock.json} : React 19, Three.js en execution ; Vite, Vitest,
Playwright en developpement) ; \textbf{images Docker} (\code{python:3.13-slim},
\code{node:22-alpine}, \code{nginx:1.27-alpine}) ; \textbf{chaine CI/CD} (actions GitHub).

Cette separation est structurante : une regression sur une dependance de developpement n'affecte
que le build, une regression sur une dependance d'execution atteint le service rendu.

\footnotesize
\begin{tabularx}{\linewidth}{@{}L{2.2cm} L{3.0cm} L{2.6cm} X@{}}
\toprule
\textbf{Type} & \textbf{Exemple} & \textbf{Mode} & \textbf{Justification} \\
\midrule
Patch \code{x.y.Z} & \code{bcrypt 5.0.0 -> 5.0.1} & \textcolor{VertOk}{\textbf{Automatique}} &
Retro-compatibilite acquise ; la CI valide la non-regression. \\
Mineure \code{x.Y.z} & \code{fastapi 0.141 -> 0.142} & \textcolor{OrangeAlerte}{\textbf{Semi-auto}} &
Auto sur \code{develop}, validation humaine avant production. \\
Majeure \code{X.y.z} & \code{React 19 -> 20} & \textcolor{RougeKo}{\textbf{Manuelle}} &
Rupture probable : branche dediee, guide de migration, tests e2e. \\
Image Docker & \code{python:3.13-slim} & \textcolor{VertOk}{\textbf{Automatique}} &
Reconstruite a chaque deploiement (correctifs systeme). \\
\bottomrule
\end{tabularx}
\normalsize

\subsection{Incident ayant motive le bornage}

\defaut{Defaut identifie --- des montees majeures non decidees}{%
\code{requirements.txt} ne declarait que des planchers (\code{>=}). Chaque installation propre et
chaque execution de CI embarquaient donc des montees \textbf{majeures} que personne n'avait
decidees ni evaluees.}

\footnotesize
\begin{tabularx}{0.86\linewidth}{@{}L{2.6cm} L{2.8cm} L{2.4cm} X@{}}
\toprule
\textbf{Dependance} & \textbf{Plancher declare} & \textbf{Installee} & \textbf{Ecart} \\
\midrule
\code{openai} & \code{>=1.51.0} & \textbf{3.3.0} & \textcolor{RougeKo}{\textbf{2 majeures}} \\
\code{rich} & \code{>=13.8.1} & \textbf{15.0.0} & \textcolor{RougeKo}{\textbf{2 majeures}} \\
\code{bcrypt} & \code{>=4.2.0} & \textbf{5.0.0} & \textcolor{OrangeAlerte}{1 majeure} \\
\code{pytest} & \code{>=8.3.2} & \textbf{9.1.1} & \textcolor{OrangeAlerte}{1 majeure} \\
\bottomrule
\end{tabularx}
\normalsize

\vspace{3pt}
Le cas sensible est \code{openai}, dont depend tout le moteur de narration. \textbf{Evaluation
d'impact} : la surface reellement appelee (\code{AsyncOpenAI}, \code{chat.completions.create} avec
\code{messages} et \code{stream}) reste disponible en 3.3.0 et les 82 tests passaient. La version
3.3.0 devient donc le \textbf{nouveau plancher valide}, plafonne a \code{<4.0.0} : la derive est
convertie en decision tracee.

\subsection{Outillage mis en place}

\textbf{1. Bornage} --- chaque dependance porte un plancher (version validee) et un plafond
(interdit la majeure). La regle « majeure = manuelle » n'est plus une consigne : elle est appliquee
par pip.

\begin{lstlisting}[style=bash]
fastapi>=0.141.1,<1.0.0     openai>=3.3.0,<4.0.0     bcrypt>=5.0.0,<6.0.0
\end{lstlisting}

\textbf{2. Detection} --- \code{.github/dependabot.yml} inspecte chaque lundi les 4 ecosystemes.
Mineures et correctifs \textbf{groupes}, majeures en \textbf{PR isolee} pour forcer une revue
individuelle.

\textbf{3. Veille} --- \code{scripts/check_updates.py} compare les versions majeures et classe les
ecarts ; code de sortie 1 si une majeure est en attente. Integre au job \code{security-audit}.

\begin{lstlisting}[style=sortie]
MONTEES MAJEURES - traitement MANUEL (3)
  npm  @testing-library/jest-dom   6.9.1 -> 7.0.1     npm  jsdom  29.1.1 -> 30.0.1
  npm  typescript                  6.0.3 -> 7.0.2
  => Ouvrir un ticket par montee, evaluer l'impact, brancher sur feature/.
MINEURES ET CORRECTIFS - traitement AUTOMATIQUE (10)  => Lot mensuel, valide par la CI.
\end{lstlisting}

\subsection{La politique observee en fonctionnement}

\begin{figure}[H]
\centering
\includegraphics[width=0.99\linewidth]{10_dependabot.png}
\caption{Pull requests reellement ouvertes par Dependabot sur le depot. Trois comportements
prevus par la configuration se lisent directement : les montees \textbf{groupees}
(\emph{frontend-runtime}, 4 mises a jour ; \emph{frontend-outils-dev}, 3 mises a jour), les montees
\textbf{majeures isolees} chacune dans sa propre PR (\code{jest-dom} 6 $\rightarrow$ 7,
\code{typescript} 6 $\rightarrow$ 7, \code{jsdom} 29 $\rightarrow$ 30), et la couverture effective
des \textbf{quatre ecosystemes} --- npm, GitHub Actions, Docker (\code{node 22-alpine}
$\rightarrow$ \code{26-alpine}) et pip.}
\end{figure}

Cette copie d'ecran n'illustre pas la politique : elle la \textbf{demontre}. Le groupement et
l'isolement des majeures ne sont pas des intentions declarees dans un fichier de configuration,
ce sont des comportements observables sur le depot. On note au passage que \code{node} passerait de
la version 22 a la version 26 : c'est exactement le type de saut qui, sans la regle d'isolement,
serait passe inapercu au milieu d'un lot.

\subsection{Frequence et evaluation d'impact}

\footnotesize
\begin{tabularx}{\linewidth}{@{}L{4.0cm} L{4.2cm} X@{}}
\toprule
\textbf{Activite} & \textbf{Frequence} & \textbf{Outillage} \\
\midrule
Audit de securite & A chaque push / PR & \code{pip-audit}, \code{npm audit} \\
Rapport de veille & A chaque push / PR & \code{scripts/check_updates.py} \\
Detection des nouvelles versions & Hebdomadaire (lundi 06h00) & \code{dependabot.yml}, 4 ecosystemes \\
Lot de mises a jour mineures & Mensuel & PR groupees, fusion apres CI verte \\
CVE elevee ou critique & Sous 48 h & Hors lot, branche \code{fix/} \\
Montees majeures & Trimestrielle & PR isolee, apres evaluation d'impact \\
\bottomrule
\end{tabularx}
\normalsize

\vspace{3pt}
Trois questions sont tranchees avant integration : la dependance est-elle \textbf{embarquee en
production} ? quelle est sa \textbf{surface d'appel} ? le correctif est-il \textbf{disponible sans
rupture} ?

\textbf{Application reelle} : l'audit npm a remonte 4 vulnerabilites « high » (\code{vite},
\code{postcss}, \code{nanoid}, \code{undici}). Toutes sont des dependances de \textbf{developpement},
absentes du bundle servi par Nginx, et leur vecteur suppose un acces au serveur de developpement.
\textbf{Decision} : lot mensuel via \code{npm audit fix}, et non correctif d'urgence sous 48 h.

\couverture
\footnotesize
\begin{tabularx}{\linewidth}{@{}L{4.0cm} X L{3.6cm}@{}}
\toprule
\textbf{Critere} & \textbf{Traitement} & \textbf{Preuve} \\
\midrule
Frequence des mises a jour & Audit et veille a chaque push, detection hebdomadaire, lot mensuel,
CVE sous 48 h, majeures trimestrielles & \code{dependabot.yml}, job \code{security-audit} \\
Perimetre logiciel & Backend Python, frontend npm, images Docker, chaine CI/CD & 4 entrees \code{updates:} \\
Type automatique / manuel & Patch auto, mineure semi-auto, majeure manuelle & Bornes \code{<N.0.0}, \code{groups} \\
\bottomrule
\end{tabularx}
\normalsize

% ===========================================================================
\section{Systeme de supervision et d'alerte}
\label{sec:supervision}

\competence{Competence C4.1.2}{Concevoir un systeme de supervision et d'alerte en determinant le
perimetre, en identifiant les indicateurs pertinents, en mettant en place des sondes et en
configurant la modalite des signalements afin de garantir une disponibilite permanente.}

\defaut{Etat initial --- une sonde qui n'interroge rien}{%
Le projet ne disposait d'aucune supervision. \code{/api/health} retournait un dictionnaire fige
affichant le \emph{chemin} de la base sans jamais l'interroger : une base corrompue, un volume non
monte ou un disque plein laissaient la sonde repondre \code{"status": "ok"} alors que toute partie
etait impossible.}

\subsection{Perimetre : ce que Prometheus et Grafana peuvent superviser}

Prometheus et Grafana supervisent le \textbf{serveur, l'API et l'infrastructure}. Ils ne mesurent
pas le rendu du jeu dans le navigateur : ce qui se passe apres la reponse HTTP --- temps de peinture,
fluidite de l'animation, erreurs JavaScript cote client --- releve d'un outillage different
(instrumentation navigateur de type RUM), hors du perimetre de ce dossier. La distinction est
importante : elle delimite ce que la supervision peut affirmer et ce qu'elle ne peut pas.

\begin{figure}[H]
\centering
\begin{tikzpicture}[
  bloc/.style={draw=AccentSec,thick,rounded corners=2pt,fill=GrisClair,align=center,
               minimum height=0.95cm,text width=2.35cm,font=\scriptsize},
  sup/.style={bloc,fill=AccentPrim!8,draw=AccentPrim},
  hors/.style={bloc,fill=white,draw=GrisBord,dashed},
  fl/.style={-{Stealth[length=2.2mm]},thick,draw=AccentSec},
  fls/.style={-{Stealth[length=2.2mm]},thick,draw=AccentPrim},
  lg/.style={font=\tiny\itshape,color=GrisTexte}]

\node[hors] (nav) {\textbf{Navigateur}\\rendu du jeu};
\node[sup, right=1.8cm of nav] (api) {\textbf{API / serveur}\\du jeu};
\node[sup, right=1.8cm of api] (bdd) {\textbf{Base de}\\\textbf{donnees}};
\node[sup, below=1.05cm of api] (prom) {\textbf{Prometheus}};
\node[sup, right=1.8cm of prom] (graf) {\textbf{Grafana}};
\node[sup, left=1.8cm of prom] (node) {\textbf{node-exporter}\\\textbf{cAdvisor}};

\draw[fl] (nav) -- node[above,lg,pos=0.5]{HTTP/SSE} (api);
\draw[fl] (api) -- node[above,lg,pos=0.5]{SQL} (bdd);
\draw[fls] (api) -- node[right,lg]{scrape} (prom);
\draw[fls] (node) -- (prom);
\draw[fls] (prom) -- (graf);

\node[lg, below=1pt of nav, text width=2.6cm, align=center] {hors perimetre\\(rendu client)};
\end{tikzpicture}
\caption{Le trait plein suit le parcours d'une requete joueur ; les fleches sombres, la chaine de
supervision. Le navigateur est le seul maillon non instrumente.}
\end{figure}

\footnotesize
\begin{tabularx}{\linewidth}{@{}L{2.7cm} X L{2.9cm}@{}}
\toprule
\textbf{Domaine} & \textbf{Ce qui est supervise} & \textbf{Par quoi} \\
\midrule
\textbf{Serveur} &
CPU, memoire, espace disque, reseau de l'hote ; consommation et redemarrages par conteneur &
\code{node-exporter}, \code{cAdvisor} \\
\addlinespace
\textbf{Backend / API} &
Temps de reponse (centiles), nombre de requetes par route, codes HTTP, exceptions non gerees &
Instrumentation applicative \\
\addlinespace
\textbf{Jeu} &
Sessions actives, flux SSE ouverts, combats engages par faction, mode de production de la narration &
Sondes metier \\
\addlinespace
\textbf{Base de donnees} &
Disponibilite verifiee par une \emph{requete reelle}, duree d'ecriture des sauvegardes, echecs
d'ecriture, volumetrie &
Sonde d'aptitude, \code{benchmark.py} \\
\addlinespace
\textbf{Erreurs} &
Taux de 4xx et 5xx, exceptions, dependance indisponible, service demarre mais inapte &
Middleware et sondes de dependance \\
\bottomrule
\end{tabularx}
\normalsize

\vspace{4pt}
\textbf{Une nuance sur la base de donnees.} Les indicateurs habituels --- nombre de connexions,
saturation du pool, requetes lentes --- supposent un SGBD \emph{serveur}. Le projet utilise
\textbf{SQLite}, moteur embarque dans le processus : il n'y a ni connexion reseau ni pool a
surveiller. Ce qui reste pertinent est donc l'\textbf{accessibilite} du fichier (verifiee par une
requete reelle a chaque collecte), la \textbf{duree d'ecriture} des sauvegardes et la
\textbf{volumetrie}. Transposer mecaniquement un tableau de bord MySQL ou PostgreSQL produirait des
panneaux vides. Une migration vers PostgreSQL justifierait en revanche l'ajout de
\code{postgres_exporter}.

\subsection{Sondes de disponibilite : vivacite contre aptitude}

\footnotesize
\begin{tabularx}{\linewidth}{@{}L{1.9cm} L{3.4cm} X L{3.4cm}@{}}
\toprule
\textbf{Sonde} & \textbf{Route} & \textbf{Finalite} & \textbf{Effet d'un echec} \\
\midrule
Vivacite & \code{GET /api/health} & Atteste que le processus repond & Redemarrage du conteneur \\
Aptitude & \code{GET /api/health/ready} & Verifie \emph{reellement} les dependances & \textbf{503} et retrait du trafic, \textbf{sans} redemarrage \\
\bottomrule
\end{tabularx}
\normalsize

\vspace{3pt}
Cette separation evite un piege d'exploitation : si la sonde de redemarrage verifiait la base, une
indisponibilite disque declencherait une \textbf{boucle de redemarrage} aggravant la panne. La sonde
d'aptitude effectue quatre controles \textbf{actifs} : requete reelle \code{SELECT count(*) FROM users},
presence et taille non nulle de \code{character_sheet.yaml} et \code{prompt_survivant.md}, ecriture
puis suppression d'un fichier temoin dans le repertoire de sauvegarde.

\begin{figure}[H]
\centering
\includegraphics[width=0.94\linewidth]{05_sonde_aptitude.png}
\caption{Sonde d'aptitude en fonctionnement : chaque dependance est controlee individuellement.}
\end{figure}

\subsection{Sondes applicatives}

La sonde structurante est \code{rpg40k_gm_generations_total{mode}}, qui distingue \code{openai},
\code{local_fallback} et \code{local_no_key} : elle rend \textbf{observable} la bascule silencieuse
vers le narrateur local, seule facon de voir une panne du fournisseur derriere un \code{200 OK}.

\footnotesize
\begin{tabularx}{\linewidth}{@{}L{6.4cm} L{1.9cm} X@{}}
\toprule
\textbf{Sonde} & \textbf{Type} & \textbf{Finalite} \\
\midrule
\code{rpg40k_gm_generations_total{mode}} & Compteur & Mode de production reel de la narration \\
\code{rpg40k_gm_generation_duration_seconds} & Histogramme & Temps de generation percu \\
\code{rpg40k_gm_context_messages} & Histogramme & Croissance de l'historique envoye au LLM \\
\code{rpg40k_http_requests_total{...}} & Compteur & Trafic et taux d'erreur par famille de code \\
\code{rpg40k_http_request_duration_seconds} & Histogramme & Centiles de latence \\
\code{rpg40k_auth_attempts_total{action,result}} & Compteur & Detection de tentatives repetees \\
\code{rpg40k_active_sessions} / \code{rpg40k_sse_streams_active} & Jauges & Capacite et charge reelle \\
\code{rpg40k_state_markers_rejected_total} & Compteur & Marqueurs d'etat ignores (\S\ref{sec:fiche}) \\
\code{rpg40k_ready} / \code{rpg40k_dependency_up} & Jauges & Etat des dependances internes \\
\bottomrule
\end{tabularx}
\normalsize

\vspace{3pt}
\textbf{Maitrise de la cardinalite} : les chemins porteurs d'identifiants sont ramenes a leur
gabarit (\code{/api/animations/tir_bolter} devient \code{/api/animations/{skill_id}}). Sans cette
normalisation, chaque identifiant creerait une serie temporelle distincte. Un test verifie que cent
identifiants ne produisent qu'un seul libelle.

\begin{figure}[H]
\centering
\includegraphics[width=0.99\linewidth,trim=0 0 0 0,clip]{04_metriques.png}
\caption{Point de collecte \texttt{/api/metrics} apres un parcours reel : version du service,
trafic par route et histogrammes de latence.}
\end{figure}

\subsection{Criteres de qualite et modalite des signalements}

\footnotesize
\begin{tabularx}{\linewidth}{@{}L{3.6cm} L{1.7cm} L{2.4cm} X@{}}
\toprule
\textbf{Indicateur} & \textbf{Objectif} & \textbf{Seuil d'alerte} & \textbf{Justification} \\
\midrule
Disponibilite API & $\geq$ 99 \%/mois & Injoignable > 1 min & Jeu solo : une coupure breve est toleree \\
Latence p95 routes de jeu & < 250 ms & > 1 s pendant 10 min & Calculs en memoire ; mesure : \textbf{1,0 ms} sur \code{/api/roll} \\
Latence \code{/api/auth/login} & < 500 ms & --- & Mesure : \textbf{$\sim$230 ms} dont \textbf{219 ms de bcrypt}. Lenteur \textbf{voulue}, a ne pas optimiser \\
Taux d'erreur 5xx & < 1 \% & > 5 \% pendant 5 min & Une erreur interrompt une partie \\
Narration servie par le LLM & 100 \% & $\geq$ 1 repli en 10 min & Tout \code{local_fallback} traduit une panne \\
Taille du contexte LLM & < 100 messages & p95 > 200 & Impact direct sur cout et latence \\
Echecs d'authentification & --- & > 10/min pendant 5 min & Signature d'une recherche exhaustive \\
\bottomrule
\end{tabularx}
\normalsize

\vspace{3pt}
Les paliers de l'histogramme (\code{0.01} a \code{10 s}) sont resserres sous la seconde, la ou se
situe la distribution reelle ; des paliers generiques auraient rendu le p95 inexploitable.

\subsection{Catalogue des 18 regles d'alerte}

\footnotesize
\begin{tabularx}{\linewidth}{@{}L{2.5cm} L{4.3cm} L{1.7cm} X@{}}
\toprule
\textbf{Groupe} & \textbf{Regle} & \textbf{Severite} & \textbf{Condition et intention} \\
\midrule
\multirow{4}{*}{\textbf{disponibilite}}
 & \code{ServiceIndisponible} & \textcolor{RougeKo}{critique} & \code{up == 0} pendant 1 min : le collecteur ne joint plus le backend \\
 & \code{ServiceNonPret} & \textcolor{RougeKo}{critique} & \code{rpg40k_ready == 0} : le processus repond mais une dependance manque \\
 & \code{DependanceIndisponible} & \textcolor{RougeKo}{critique} & Identifie \emph{laquelle} des quatre dependances est tombee \\
 & \code{SondeExterneEnEchec} & \textcolor{RougeKo}{critique} & Sonde boite noire : couvre aussi une panne du reverse proxy \\
\midrule
\multirow{5}{*}{\parbox{2.4cm}{\raggedright\textbf{ressources serveur}}}
 & \code{EspaceDisqueFaible} & \textcolor{OrangeAlerte}{majeure} & < 15 \% d'espace libre : les sauvegardes echoueront avant tout signal applicatif \\
 & \code{EspaceDisqueCritique} & \textcolor{RougeKo}{critique} & < 5 \% : perte de donnees imminente \\
 & \code{MemoireHoteSaturee} & \textcolor{OrangeAlerte}{majeure} & > 90 \% pendant 10 min : risque d'arret par le noyau \\
 & \code{ChargeProcesseurElevee} & mineure & > 85 \% pendant 15 min : sur un service a faible trafic, signale plutot une boucle \\
 & \code{ConteneurRedemarrages} & \textcolor{OrangeAlerte}{majeure} & Redemarrage inattendu d'un conteneur \\
\midrule
\multirow{3}{*}{\parbox{2.4cm}{\raggedright\textbf{qualite de service}}}
 & \code{TauxErreurServeurEleve} & \textcolor{OrangeAlerte}{majeure} & > 5 \% de 5xx sur 5 min \\
 & \code{LatenceApiDegradee} & \textcolor{OrangeAlerte}{majeure} & p95 > 1 s sur 10 min, hors routes de streaming \\
 & \code{ExceptionsNonGerees} & \textcolor{OrangeAlerte}{majeure} & Toute exception remontee jusqu'au middleware \\
\midrule
\multirow{3}{*}{\parbox{2.4cm}{\raggedright\textbf{dependance LLM}}}
 & \code{ModeDegradeNarration} & \textcolor{OrangeAlerte}{majeure} & \textbf{Regle structurante} : tout repli sur le narrateur local \\
 & \code{LatenceLlmElevee} & mineure & p95 de generation > 20 s \\
 & \code{ContexteConversationExcessif} & mineure & p95 > 200 messages : derive de cout et de latence \\
\midrule
\textbf{securite} & \code{EchecsAuthentificationRepetes} & \textcolor{OrangeAlerte}{majeure} & > 10 echecs/min : recherche exhaustive probable \\
\midrule
\multirow{2}{*}{\textbf{capacite}}
 & \code{SessionsMemoireCroissantes} & mineure & > 200 sessions residentes pendant 15 min \\
 & \code{EchecSauvegarde} & \textcolor{RougeKo}{critique} & Toute ecriture de sauvegarde en echec : perte de progression \\
\bottomrule
\end{tabularx}
\normalsize

\vspace{4pt}
Les fenetres \code{for:} sont volontairement superieures a l'intervalle de collecte, afin qu'un pic
isole ne declenche pas d'alerte.

\subsection{Modalite des signalements}

Le routage Alertmanager traduit la severite en canal \emph{et} en cadence.

\footnotesize
\begin{tabularx}{\linewidth}{@{}L{2.1cm} L{4.0cm} c c X@{}}
\toprule
\textbf{Severite} & \textbf{Destinataire} & \textbf{Groupe} & \textbf{Rappel} & \textbf{Intention} \\
\midrule
\textcolor{RougeKo}{\textbf{Critique}} & Astreinte, notification immediate & 10 s & 1 h & Service indisponible ou donnees menacees \\
\textcolor{OrangeAlerte}{\textbf{Majeure}} & Equipe, heures ouvrees & 1 min & 4 h & Degradation perceptible par le joueur \\
\textcolor{AccentSec}{\textbf{Mineure}} & Consignation seule & 30 s & 24 h & Derive a surveiller, revue hebdomadaire \\
\bottomrule
\end{tabularx}
\normalsize

\vspace{4pt}
Deux mecanismes limitent le bruit, condition d'une supervision \emph{reellement} exploitee. Le
\textbf{groupement} (\code{group_by: [alertname, severite]}) fait qu'une panne declenchant dix
regles produit une notification par famille et non dix. L'\textbf{inhibition} supprime les alertes
de latence et de taux d'erreur lorsque \code{ServiceIndisponible} est actif, puisqu'elles en
decoulent mecaniquement. Les URL de destination ne sont jamais versionnees : elles proviennent des
variables \code{ALERT_WEBHOOK_CRITIQUE} et \code{ALERT_WEBHOOK_EQUIPE}.

\subsection{Architecture de la chaine}

\begin{figure}[H]
\centering
\begin{tikzpicture}[
  bloc/.style={draw=AccentSec,thick,rounded corners=2pt,fill=GrisClair,align=center,
               minimum height=1.0cm,text width=2.5cm,font=\scriptsize},
  alt/.style={bloc,fill=white},
  fl/.style={-{Stealth[length=2.2mm]},thick,draw=AccentSec},
  lg/.style={font=\tiny\itshape,color=GrisTexte}]
\node[bloc] (back) {\textbf{Backend}\\\texttt{/api/metrics}\\\texttt{/health/ready}};
\node[bloc, right=1.5cm of back] (prom) {\textbf{Prometheus}\\collecte 15 s\\retention 30 j};
\node[bloc, right=1.5cm of prom] (alert) {\textbf{Alertmanager}\\routage, groupage\\inhibition};
\node[alt, below=0.8cm of back] (nginx) {\textbf{Nginx + React}};
\node[alt, below=0.8cm of prom] (black) {\textbf{Blackbox}\\sonde externe};
\node[bloc, below=0.8cm of alert] (graf) {\textbf{Grafana}\\13 panneaux};
\draw[fl] (back) -- node[above,lg]{scrape 15 s} (prom);
\draw[fl] (prom) -- node[above,lg]{18 regles} (alert);
\draw[fl] (prom) -- (graf);
\draw[fl] (black) -- (prom);
\draw[fl] (nginx) -- (black);
\draw[fl] (back) -- (nginx);
\end{tikzpicture}
\caption{Collecte applicative, sonde externe « boite noire », evaluation des regles et
visualisation. Prometheus se supervise lui-meme : sans cela, l'arret du collecteur produirait un
silence indistinguable d'une absence d'incident.}
\end{figure}

La pile s'active en surcouche, ce qui laisse la pile applicative deployable seule. Prometheus,
Alertmanager et Grafana n'exposent leurs ports que sur \code{127.0.0.1} : ils n'embarquent pas
d'authentification suffisante pour une exposition publique, et l'acces distant passe par un tunnel
SSH.

\begin{lstlisting}[style=bash]
docker compose -f docker-compose.yml -f docker-compose.monitoring.yml up -d
\end{lstlisting}

\subsection{Validation de la pile en fonctionnement}

La chaine a ete executee de bout en bout pour verifier qu'elle ne se limite pas a des fichiers de
configuration valides. Prometheus a collecte le backend reel, evalue les 18 regles, et Grafana a
affiche le tableau de bord \emph{provisionne depuis le depot}.

Premiere verification, avec \code{promtool}, l'outil officiel de Prometheus --- il controle la
syntaxe \emph{et} les expressions PromQL, ce qu'une simple validation YAML ne fait pas :

\begin{lstlisting}[style=sortie]
$ promtool check rules alert_rules.yml      $ promtool check config prometheus.yml
  SUCCESS: 18 rules found                     SUCCESS: prometheus.yml is valid

$ curl -s localhost:9090/api/v1/targets
  prometheus       http://127.0.0.1:9090/metrics       UP
  rpg40k-backend   http://127.0.0.1:8000/api/metrics   UP
\end{lstlisting}

Une panne du fournisseur LLM a ensuite ete simulee (cle API invalide) pendant que Prometheus
collectait. L'alerte structurante est passee en etat \textbf{firing} :

\begin{figure}[H]
\centering
\includegraphics[width=0.99\linewidth]{11_alertes.png}
\caption{Regles evaluees par Prometheus. Le groupe \texttt{dependance\_llm} affiche
\textbf{FIRING (1)} : \texttt{ModeDegradeNarration} s'est declenchee parce que la narration etait
servie par le narrateur local. L'application repondait pourtant \texttt{200 OK} --- c'est exactement
la degradation qu'aucune supervision fondee sur les codes HTTP n'aurait vue.}
\end{figure}

\begin{figure}[H]
\centering
\includegraphics[width=0.99\linewidth]{14_grafana.png}
\caption{Tableau de bord d'exploitation, alimente par des donnees reelles. Le panneau
\emph{Mode de production de la narration} distingue \texttt{local\_fallback} de \texttt{local\_no\_key} ;
le taux d'erreur 5xx affiche \emph{No data}, ce qui est exact : aucune erreur serveur n'est
survenue pendant la mesure.}
\end{figure}

\subsection{Tableau de bord infrastructure et jeu}

Un second tableau de bord regroupe les indicateurs d'exploitation courants : ressources de l'hote,
activite de jeu et sante de l'API sur une meme vue.

\begin{figure}[H]
\centering
\includegraphics[width=0.99\linewidth]{15_infrastructure.png}
\caption{Huit indicateurs alimentes par des mesures reelles : charge processeur, memoire utilisee,
debit reseau, joueurs connectes, latence p95, erreurs cumulees, uptime du processus et debit de
requetes. En bas, l'evolution processeur/memoire et l'espace disque par volume, avec les seuils de
15 \% et 5 \% des regles \texttt{EspaceDisque}.}
\end{figure}

\footnotesize
\begin{tabularx}{\linewidth}{@{}L{3.1cm} L{5.6cm} X@{}}
\toprule
\textbf{Indicateur} & \textbf{Source} & \textbf{Interet} \\
\midrule
CPU & \texttt{node\_cpu\_seconds\_total} & Saturation de l'hote \\
RAM utilisee & \texttt{node\_memory\_*} & Risque d'arret par le noyau \\
Reseau & \texttt{node\_network\_*\_bytes\_total} & Volume echange, detection d'anomalie \\
Joueurs connectes & \texttt{rpg40k\_active\_sessions} & Charge reelle et derive memoire \\
Latence p95 & \texttt{rpg40k\_http\_request\_duration\_seconds} & Experience percue, hors streaming \\
Erreurs (1 h) & \texttt{rpg40k\_http\_requests\_total} & Volume de 4xx et 5xx \\
Uptime & \texttt{rpg40k\_process\_start\_time\_seconds} & Detecte les redemarrages silencieux \\
Requetes / s & \texttt{rpg40k\_http\_requests\_total} & Debit traite par l'API \\
\bottomrule
\end{tabularx}
\normalsize

\vspace{4pt}
La sonde d'uptime a ete ajoutee pour ce tableau de bord : \code{prometheus_client} n'expose
\code{process_start_time_seconds} que sous Linux, ce qui aurait laisse la tuile vide sur une autre
plateforme. Les panneaux de ressources sont ecrits avec l'operateur \code{or} de PromQL
(\code{node_...} \emph{ou} \code{windows_...}), de sorte que le meme fichier fonctionne en
production Linux comme en validation locale.

\reserve{%
\textbf{Portee de cette validation.} Docker n'etant pas disponible sur le poste de travail, la
pile a ete executee a partir des \textbf{binaires autonomes} de Prometheus 3.1.0 et Grafana 11.5.1,
avec les cibles adressees en \texttt{127.0.0.1} au lieu des noms de services Compose. Les
\textbf{regles d'alerte et le tableau de bord sont ceux du depot, sans modification} ; seules les
adresses des cibles different. \texttt{node-exporter}, \texttt{cAdvisor}, \texttt{blackbox-exporter} et
Alertmanager restent donc configures mais non executes.}

\subsection{Couverture des modes de defaillance}

\footnotesize
\begin{tabularx}{\linewidth}{@{}L{4.4cm} L{4.2cm} X@{}}
\toprule
\textbf{Mode de defaillance} & \textbf{Detection} & \textbf{Reaction} \\
\midrule
Processus arrete ou fige & \code{up == 0} (1 min) & Redemarrage automatique, alerte critique \\
Dependance en panne & \code{rpg40k_ready == 0} (2 min) & Retrait du trafic (503), \textbf{pas} de redemarrage \\
Panne du proxy ou du frontend & \code{probe_success == 0} & Alerte critique meme si le backend est sain \\
Service disponible mais \textbf{degrade} & \code{mode="local_fallback"} & Alerte majeure : degradation rendue visible \\
\bottomrule
\end{tabularx}
\normalsize

\vspace{3pt}
Le dernier cas est celui que la sonde initiale ne pouvait pas voir, et c'est le plus probable en
exploitation : cle expiree, quota depasse, incident du fournisseur. \textbf{20 tests} couvrent le
dispositif (82 $\rightarrow$ 102) et verifient que la sonde d'aptitude detecte \emph{reellement}
une base corrompue et un fichier de jeu absent.

\couverture
\footnotesize
\begin{tabularx}{\linewidth}{@{}L{4.4cm} X@{}}
\toprule
\textbf{Critere} & \textbf{Traitement} \\
\midrule
Systeme adapte a la typologie & Quatre caracteristiques du logiciel traduites en choix de sondes et de seuils (\S1) \\
Sondes et finalite explicitees & Disponibilite, applicatives, LLM, securite, capacite : type et finalite pour chacune \\
Criteres de qualite adaptes & Sept indicateurs, seuils justifies par des mesures reelles (1,0 ms ; 230 ms dont 219 de bcrypt) \\
Surveillance de la disponibilite & Quatre modes de defaillance couverts, dont la degradation silencieuse \\
\bottomrule
\end{tabularx}
\normalsize

% ===========================================================================
\section{Collecte et consignation des anomalies}
\label{sec:collecte}

\competence{Competence C4.2.1}{Consigner les anomalies detectees en elaborant un processus de
collecte et consignation, en utilisant des outils de collecte et en y integrant toutes les
informations pertinentes, afin de determiner le correctif a mettre en place.}

\subsection{Adaptation a la typologie du logiciel}

\textbf{1. Sorties non deterministes.} La narration provenant d'un LLM, une anomalie declenchee par
une formulation particuliere est \textbf{intermittente et non rejouable a l'identique}. La fiche
impose donc \emph{Frequence} et \emph{Horodatage}, seuls moyens de retrouver le texte incrimine
dans les journaux.

\textbf{2. Degradation sans erreur visible.} Une part des anomalies ne se manifeste \textbf{jamais}
par un code d'erreur : elles ne sont detectables que par les sondes du \S\ref{sec:supervision}.
Un formulaire distinct existe pour ces incidents, qui n'ont pas de signalement utilisateur.

\textbf{3. Sauvegarde par scene.} Une anomalie en fin de traitement peut faire perdre une scene
entiere : la gravite est evaluee a l'aune de la \textbf{perte de donnees}, non du seul blocage.

\subsection{Sources, outils et cycle de vie}

\footnotesize
\begin{tabularx}{\linewidth}{@{}L{3.2cm} X L{2.6cm}@{}}
\toprule
\textbf{Source} & \textbf{Outil} & \textbf{Delai} \\
\midrule
Supervision & 18 regles Prometheus / Alertmanager & Temps reel \\
Journaux applicatifs & Journalisation structuree & Temps reel \\
Integration continue & GitHub Actions (5 jobs) & A chaque push \\
Audit de dependances & \code{pip-audit}, \code{npm audit} & A chaque push \\
Signalement joueur & Formulaire GitHub Issues & Variable \\
Recette manuelle & Plan de recette & Avant production \\
\bottomrule
\end{tabularx}
\normalsize

\vspace{3pt}
Les quatre premieres sources sont \textbf{automatiques} : point decisif pour un logiciel a faible
base d'utilisateurs, ou l'on ne peut compter sur le volume de signalements.

La collecte est centralisee dans GitHub Issues avec trois formulaires structures, la saisie libre
etant desactivee (\code{blank_issues_enabled: false}) : \code{01_anomalie.yml} (\textbf{10 champs
obligatoires}), \code{02_incident_supervision.yml} (5), \code{03_retour_utilisateur.yml} (4). Le
formulaire contraint repond a un constat pratique : une description en prose omet presque toujours
la version, l'environnement ou les etapes exactes, imposant un aller-retour avant analyse. Les dix
champs obligatoires sont exactement ceux \textbf{sans lesquels le bogue n'est pas reproductible}.

\begin{figure}[H]
\centering
\begin{tikzpicture}[
  bloc/.style={draw=AccentSec,thick,rounded corners=2pt,fill=GrisClair,align=center,
               minimum height=1.15cm,text width=2.55cm,font=\scriptsize},
  fl/.style={-{Stealth[length=2.2mm]},thick,draw=AccentPrim},
  et/.style={font=\scriptsize\itshape,color=AccentPrim}]
\node[bloc] (d) {\textbf{Detection}\\supervision, journaux, CI, joueur};
\node[bloc, right=0.95cm of d] (q) {\textbf{Qualification}\\reproduction, gravite, priorisation};
\node[bloc, right=0.95cm of q] (t) {\textbf{Traitement}\\branche \texttt{fix/}, correctif + test};
\node[bloc, right=0.95cm of t] (c) {\textbf{Cloture}\\test de non-regression, fermeture};
\draw[fl] (d) -- node[above,font=\tiny\itshape]{fiche} (q);
\draw[fl] (q) -- (t); \draw[fl] (t) -- (c);
\node[et, below=1pt of d] {a-qualifier}; \node[et, below=1pt of q] {confirmee};
\node[et, below=1pt of t] {en-cours}; \node[et, below=1pt of c] {corrigee};
\end{tikzpicture}
\caption{Cycle de vie d'une anomalie ; les etats correspondent aux etiquettes GitHub et au kanban.}
\end{figure}

\textbf{Grille de gravite} --- \textcolor{RougeKo}{\textbf{Bloquante}} (jeu inutilisable,
immediat) ; \textcolor{OrangeAlerte}{\textbf{Majeure}} (fonctionnalite inutilisable \emph{ou perte
de donnees}, 48 h) ; \textbf{Mineure} (gene sans perte, lot mensuel) ; \textbf{Cosmetique}
(affichage, opportuniste). La perte de donnees est classee majeure \emph{meme sans blocage} : dans
un jeu a sauvegarde par scene, perdre une scene annule plusieurs minutes de jeu.

% ===========================================================================
\section{Fiche de consignation --- ANO-2026-001}
\label{sec:fiche}

\competence{Competence C4.2.1 (livrable 2)}{La presentation d'une fiche de consignation d'une
anomalie rencontree au cours du projet.}

\footnotesize
\begin{tabularx}{\linewidth}{@{}L{2.6cm} X L{2.4cm} L{3.2cm}@{}}
\toprule
\textbf{Identifiant} & \code{ANO-2026-001} & \textbf{Gravite} & \textcolor{OrangeAlerte}{\textbf{Majeure}} (perte de donnees) \\
\textbf{Perimetre} & Narration / maitre du jeu & \textbf{Version} & \code{v1.0.0-rncp} \\
\textbf{Frequence} & Intermittente, depend du texte du LLM & \textbf{Etat} & \textcolor{VertOk}{\textbf{Corrigee}} en 1.3.0 \\
\bottomrule
\end{tabularx}
\normalsize

\subsection{Description et reproduction}

A l'issue de certaines scenes, le joueur recoit integralement la narration puis l'interface reste
bloquee sur l'indicateur d'attente, sans message d'erreur. En rechargeant, \textbf{la scene qui
vient d'etre jouee a disparu}. Le phenomene est intermittent et independant de l'action du joueur,
ce qui orientait a tort vers un probleme reseau.

Le declencheur etant une sortie du LLM, la reproduction fiable passe par la substitution du
narrateur : demarrer une partie, puis faire produire une narration dont la valeur de ressource
n'est \textbf{pas numerique}, et observer le flux SSE de \code{POST /api/chat}.

\begin{lstlisting}[style=sortie]
Karimus partage ses dernieres provisions.  [ETAT: rations=aucune, stress=+1]

--- NOMINAL  (rations=-1) ---            --- DEFAILLANT (rations=aucune) ---
  statut HTTP        : 200                 !! ValueError: invalid literal for
  evenements 'token' : 5                      int() with base 10: 'aucune'
  evenements 'done'  : 1                   -> aucun evenement 'done' emis
\end{lstlisting}

\subsection{Analyse technique}

\textbf{Cause racine} --- dans \code{src/state.py}, \code{apply_updates_from_text()} convertit la
valeur d'une ressource sans proteger la conversion. \textbf{Asymetrie revelatrice} : la branche
voisine, qui traite \code{stress} et \code{corruption}, intercepte pourtant \emph{exactement} le
meme cas. Le defaut n'est donc pas un oubli de conception mais une \textbf{protection appliquee a
une branche et omise sur l'autre}.

\begin{lstlisting}[style=python]
# branche resources : conversion nue -> ValueError
if value.startswith(('+','-')): delta = int(value)
else:                           self.set_resource(key, int(value))

# branche tracks : le meme cas est deja intercepte
try:                self.tracks[key] = int(value)
except ValueError:  self.tracks[key] = value
\end{lstlisting}

\textbf{Mecanisme de l'impact} --- l'appel se situe dans le generateur SSE ; l'ordre des
instructions explique les deux symptomes simultanes.

\begin{lstlisting}[style=python]
546  session.messages.append({"role":"assistant","content": full_text})
548  changes = session.character.apply_updates_from_text(full_text)  # exception ici
550  session.save()                                                  # jamais atteint
551  yield json.dumps({"type":"done", ...})                          # jamais emis
\end{lstlisting}

L'exception survient \textbf{apres} l'envoi de la narration (le joueur a lu la scene), \textbf{avant}
\code{session.save()} (progression perdue) et \textbf{avant} l'evenement \code{done} (client bloque).
Le defaut avait echappe aux tests parce que la suite ne couvrait ce parseur qu'avec des valeurs
numeriques valides.

\subsection{Preconisations et verification}

\footnotesize
\begin{tabularx}{\linewidth}{@{}L{4.2cm} X c@{}}
\toprule
\textbf{Option} & \textbf{Analyse} & \textbf{Retenue} \\
\midrule
A. Renforcer le prompt systeme & Ne supprime pas le risque : la sortie d'un LLM ne se garantit pas par une consigne & \textcolor{RougeKo}{Non} \\
B. Intercepter la \code{ValueError} dans le parseur & Traite la cause racine, aligne \code{resources} sur \code{tracks} & \textcolor{VertOk}{\textbf{Oui}} \\
C. Proteger l'appel dans le generateur SSE & Ne traite pas la cause, mais garantit qu'aucune future erreur ne coute une sauvegarde & \textcolor{VertOk}{\textbf{Oui}} \\
\bottomrule
\end{tabularx}
\normalsize

\vspace{3pt}
B et C sont retenues ensemble, en \textbf{defense en profondeur} : B corrige le defaut identifie,
C protege contre les defauts de meme famille non encore decouverts. Une quatrieme mesure rend le
phenomene \textbf{mesurable} : les sondes \code{rpg40k_state_markers_rejected_total} et
\code{rpg40k_state_parse_failures_total}. Un marqueur ignore reste en effet une anomalie mineure
(la ressource n'est pas mise a jour alors que la narration le laisse entendre) : la sonde permet
d'en mesurer la frequence residuelle.

\textbf{18 tests de non-regression} (dont 8 valeurs textuelles en parametre), portant la suite de
102 a \textbf{120}. Apres correctif, le meme parcours renvoie \code{200}, 6 evenements \code{token},
1 evenement \code{done}, et la sonde de marqueurs rejetes vaut \code{1.0}.

\couverture
\footnotesize
\begin{tabularx}{\linewidth}{@{}L{4.4cm} X@{}}
\toprule
\textbf{Critere} & \textbf{Traitement} \\
\midrule
Processus structure et adapte & Six sources dont quatre automatiques, trois formulaires, grille fondee sur la perte de donnees, cycle en quatre etats \\
Informations permettant de reproduire & Etapes, marqueur declencheur exact, commande de reproduction automatisee, releves avant/apres \\
Analyse et preconisations explicitees & Cause racine, asymetrie des branches, impact ligne a ligne ; trois options comparees, deux retenues, 18 tests \\
\bottomrule
\end{tabularx}
\normalsize

% ===========================================================================
\section{Traitement de l'anomalie par la chaine CI/CD}
\label{sec:cicd}

\competence{Competence C4.2.2}{Creer et deployer un correctif en respectant le processus
d'integration et de deploiement continu afin de resoudre l'anomalie.}

\subsection{Chaine mobilisee}

\footnotesize
\begin{tabularx}{\linewidth}{@{}c L{4.0cm} L{3.4cm} X@{}}
\toprule
\textbf{\#} & \textbf{Etape} & \textbf{Declencheur} & \textbf{Controle exerce} \\
\midrule
1 & Branche de correction & Manuel & Isolation du correctif \\
2 & Verification locale & Manuel & Boucle courte avant push \\
3 & Integration continue & Push / pull request & 5 jobs GitHub Actions \\
4 & Revue & Pull request & Relecture du correctif et des tests \\
5 & Fusion vers \code{main} & Manuel apres CI verte & Seul un code teste atteint \code{main} \\
6 & Deploiement & Automatique sur CI verte & Reconstruction et redeploiement Docker \\
7 & Validation post-deploiement & Automatique & 5 controles fonctionnels \\
8 & Retour arriere & Automatique si echec & Restauration du commit precedent \\
\bottomrule
\end{tabularx}
\normalsize

\subsection{Decoupage en commits}

Le correctif est decoupe pour que \textbf{chaque commit reste independamment testable}, condition
pour qu'une bissection (\code{git bisect}) demeure exploitable et qu'un retour arriere partiel soit
possible.

\begin{lstlisting}[style=sortie]
9f98298  test(animations): rendre deterministe le test de fin d'animation (ANO-2026-003)
0098252  fix(sauvegarde): persister la fiche de personnage (ANO-2026-002)
ee546a3  feat(version): source unique de verite et journal des versions 1.3.0
5cafcbd  feat(perf): outiller la mesure des indicateurs et la collecte des retours
8da4dca  ci(deploy): valider le deploiement par test de fumee et rollback automatique
26df2a1  fix(narration): ne plus perdre la progression sur un marqueur d'etat invalide
b335510  feat(supervision): instrumenter l'application et outiller l'alerte
e9e52d5  chore(deps): borner les versions et automatiser la veille des dependances
\end{lstlisting}

L'ordre n'est pas arbitraire : le commit de supervision \textbf{precede} le correctif, car
celui-ci s'appuie sur les sondes qu'il introduit. Le message suit la convention \emph{Conventional
Commits} deja employee dans le depot et reference explicitement l'anomalie, ce qui relie le commit
a sa fiche de consignation et permet d'en retrouver le contexte des mois plus tard.

\begin{lstlisting}[style=sortie]
fix(narration): ne plus perdre la progression sur un marqueur d'etat invalide
...
Fixes: ANO-2026-001
Refs: C4.2.1
\end{lstlisting}

\subsection{Preuve de resolution}

La validite est etablie en executant la \textbf{meme suite de tests} avant puis apres le correctif.
Ce double passage est la demonstration recherchee : les tests \textbf{echouent effectivement} sur le
code d'origine, ce qui prouve qu'ils decrivent l'anomalie et non un comportement quelconque. Les
deux tests qui passaient deja sont ceux du comportement nominal, que le correctif ne devait pas
modifier.

\footnotesize
\begin{tabularx}{0.92\linewidth}{@{}L{5.4cm} L{4.0cm} X@{}}
\toprule
\textbf{Commit} & \textbf{Resultat} & \\
\midrule
\code{b335510} --- precedant le correctif & \textcolor{RougeKo}{\textbf{16 failed}}, 2 passed & \\
\code{26df2a1} --- commit du correctif & \textcolor{VertOk}{\textbf{18 passed}} & \\
\bottomrule
\end{tabularx}
\normalsize

\subsection{Validation post-deploiement et retour arriere}

\defaut{Defaut identifie --- un deploiement valide sur la mauvaise sonde}{%
Le deploiement se terminait sur un simple \code{curl} de \code{/api/health}, la \textbf{sonde de
vivacite}, qui repond \code{ok} des le demarrage du processus : une version inapte au trafic etait
declaree deployee avec succes. \textbf{Aucun retour arriere n'existait.}}

Le workflow memorise desormais le commit en service, execute \code{scripts/smoke_test.sh}, et
restaure la version precedente en cas d'echec.

\begin{lstlisting}[style=bash]
# Version actuellement en service, conservee pour un retour arriere.
VERSION_PRECEDENTE="$(git -C "$APP_DIR" rev-parse HEAD)"

docker compose -p rpg40k up -d --build

if ! ./scripts/smoke_test.sh "http://127.0.0.1:$RPG40K_HTTP_PORT"; then
  echo "::error::Test de fumee en echec - retour a $VERSION_PRECEDENTE"
  git -C "$APP_DIR" checkout --force "$VERSION_PRECEDENTE"
  docker compose -p rpg40k up -d --build
  exit 1
fi
\end{lstlisting}

Le test exerce cinq controles, choisis pour distinguer un service \emph{demarre} d'un service
\emph{operationnel}.

\footnotesize
\begin{tabularx}{\linewidth}{@{}c L{5.0cm} X@{}}
\toprule
\textbf{\#} & \textbf{Controle} & \textbf{Ce qu'il detecte} \\
\midrule
1 & Disponibilite (30 tentatives) & Conteneur non demarre ou build echoue \\
2 & \textbf{Sonde d'aptitude} \code{/api/health/ready} & Base injoignable, volume non monte, disque plein \\
3 & Exposition de \code{/api/metrics} & Version deployee sans supervision --- angle mort \\
4 & \code{/api/state} sans jeton renvoie \code{401} & Regression de configuration exposant les routes de jeu \\
5 & Livraison de l'interface & Panne du reverse proxy Nginx \\
\bottomrule
\end{tabularx}
\normalsize

\begin{lstlisting}[style=sortie]
1. Attente de la disponibilite     [OK]  service joignable apres 1 tentative(s)
2. Sonde d'aptitude                [OK]  toutes les dependances sont disponibles
3. Point de collecte des metriques [OK]  metriques exposees
4. Protection des routes de jeu    [OK]  acces non authentifie refuse (401)
5. Livraison de l'interface        [ECHEC] interface non servie par le reverse proxy
\end{lstlisting}

Le cinquieme controle echoue \textbf{volontairement} ici : le script visait le backend seul
(port 8000), qui ne sert pas l'interface ; en production la cible est le port du reverse proxy. Ce
releve demontre que le script \textbf{detecte reellement une chaine incomplete} plutot que de
valider systematiquement.

\textbf{Garde-fou de securite ajoute au passage} : le deploiement creait \code{.env} par copie de
\code{.env.example}, ce qui signait les jetons de production avec le \code{JWT_SECRET} d'exemple
\emph{public dans le depot}. Le deploiement echoue desormais si le secret est absent, trop court ou
laisse a sa valeur d'exemple.

\begin{figure}[H]
\centering
\includegraphics[width=0.99\linewidth]{06_ci_github.png}
\caption{Executions reelles de la chaine. La CI \#86 est verte sur \texttt{main} ; Dependabot a
ouvert 11 pull requests, dont la montee \emph{majeure} \code{jest-dom} 6 $\rightarrow$ 7 en PR
isolee, conformement a la politique du \S\ref{sec:deps}.}
\end{figure}

\couverture
\footnotesize
\begin{tabularx}{\linewidth}{@{}L{4.4cm} X@{}}
\toprule
\textbf{Critere} & \textbf{Traitement} \\
\midrule
Le traitement tire profit de la chaine CI/CD & Branche dediee, commits atomiques referencant l'anomalie, validation par 5 jobs, deploiement conditionne a une CI verte, test de fumee, rollback automatique \\
Le correctif est decrit et resout l'anomalie & Deux niveaux de correctif code a l'appui ; memes tests avant (16 echecs) et apres (18 succes) ; releve avant/apres sur l'application \\
\bottomrule
\end{tabularx}
\normalsize

% ===========================================================================
\section{Recommandations argumentees d'amelioration}
\label{sec:reco}

\competence{Competence C4.3.1}{Proposer des axes d'amelioration en prenant en compte les
indicateurs de performance et en analysant les retours utilisateurs afin de maintenir et renforcer
l'attractivite du logiciel.}

\reserve{\textbf{Limite assumee.} Le projet n'a pas de base d'utilisateurs constituee : aucun retour
reel n'a pu etre analyse. Le canal de collecte (\code{03_retour_utilisateur.yml}) a donc ete
\textbf{mis en place} au titre de cette competence, et les recommandations s'appuient sur les
indicateurs mesures et les observations de recette. Presenter des retours fictifs aurait prive
l'analyse de toute valeur probante.}

\subsection{Indicateurs mesures}

Releve reproductible par \code{python scripts/benchmark.py} :

\begin{lstlisting}[style=sortie]
LATENCE (ms)      p50     p95      POIDS LIVRE AU NAVIGATEUR
/api/health      2.59     3.1        VoxelEngine.js    481.4 Ko   <- 64 % du total
/api/roll         3.2    3.82        index.js          245.7 Ko
/api/state        9.97  11.02        TOTAL JS          756.3 Ko
/api/auth/login 248.56 287.24
                                    BASE : users 87, session_events 141, index : aucun
CONTEXTE LLM   scene 1 : 1209 tokens | scene 10 : 3369 | scene 30 : 8169 (cumul 140692)
\end{lstlisting}

\textbf{Lecture.} Les routes de jeu (3 a 10 ms) ne sont pas un facteur limitant : aucune
optimisation n'est justifiee de ce cote. Les deux points de tension sont le \textbf{poids livre au
navigateur} et la \textbf{croissance du contexte LLM}. La latence du login est \textbf{conforme a
l'attendu} (219 ms de \code{bcrypt.checkpw}) et ne doit pas etre reduite.

\subsection{R1 --- Charger le moteur 3D a la demande}

\footnotesize
\begin{tabularx}{\linewidth}{@{}L{2.5cm} X@{}}
\toprule
\textbf{Indicateur} & \code{VoxelEngine.js} pese \textbf{481 Ko} sur 756 Ko de JavaScript, soit \textbf{64 \% du poids livre} \\
\textbf{Constat} & Le jeu est avant tout \textbf{textuel}. Un joueur qui ne declenche jamais de combat 3D telecharge pourtant tout le moteur voxel et Three.js \\
\textbf{Cout / delai} & \textcolor{VertOk}{Faible} --- 0,5 a 1 jour ; immediat, sans rupture fonctionnelle \\
\textbf{Gain attendu} & \textbf{$-$64 \% de JavaScript au premier chargement} ; premier rendu nettement plus rapide, en particulier sur connexion mobile \\
\textbf{Risque} & Faible : latence au premier combat 3D, absorbable par un indicateur de chargement \\
\bottomrule
\end{tabularx}
\normalsize

\vspace{3pt}
Vite produit \emph{deja} des fragments separes (\code{VoxelEngine-*.js}, \code{Combat3DDemo-*.js}) :
le decoupage existe, seuls les imports restent statiques. Il suffit de les convertir en imports
dynamiques (\code{React.lazy} / \code{import()}). Le temps d'affichage initial etant le premier
facteur d'abandon d'une application web, c'est le meilleur rapport gain/effort du lot.

\subsection{R2 --- Borner l'historique envoye au LLM}

\footnotesize
\begin{tabularx}{\linewidth}{@{}L{2.5cm} X@{}}
\toprule
\textbf{Indicateur} & Contexte de \textbf{1 209} tokens (scene 1) a \textbf{8 169} (scene 30) ; cumul envoye \textbf{140 692} tokens \\
\textbf{Constat} & \code{session.messages} n'est jamais tronque : chaque scene renvoie tout l'historique \\
\textbf{Cout / delai} & \textcolor{OrangeAlerte}{Faible a moyen} --- 1 a 2 jours ; court terme \\
\textbf{Gain attendu} & Croissance \textbf{lineaire au lieu de quadratique} ; latence stabilisee ; plafond de contexte jamais atteint \\
\textbf{Risque} & \textcolor{OrangeAlerte}{Moyen} --- une troncature naive fait « oublier » au maitre du jeu les evenements anciens \\
\bottomrule
\end{tabularx}
\normalsize

\vspace{3pt}
\textbf{Argument de cout, presente honnetement.} Le gain financier est faible en valeur absolue :
environ \textbf{0,021 USD} par partie de 30 scenes contre \textbf{0,015 USD} avec une fenetre
glissante de 10 messages. Ce n'est \textbf{pas} l'argument principal, et le presenter comme tel
serait trompeur. Les deux vraies justifications sont ailleurs : la \textbf{latence}, qui croit avec
le contexte et rompt le rythme de jeu au-dela de quelques dizaines de scenes ; et le \textbf{plafond
de contexte} du modele, dont le depassement provoquerait une erreur, donc aujourd'hui une bascule
en mode degrade.

\textbf{Mise en oeuvre recommandee} : conserver le prompt systeme, les $N$ derniers echanges et un
\textbf{resume condense} des scenes anterieures, regenere tous les 10 tours --- approche qui
preserve la continuite narrative, precisement ce que le joueur percoit. La sonde
\code{rpg40k_gm_context_messages} mesurera le gain reel.

\subsection{R3 a R5 --- Capacite, reproductibilite et securite}

\footnotesize
\begin{tabularx}{\linewidth}{@{}L{1.1cm} L{3.4cm} X L{1.3cm}@{}}
\toprule
\textbf{\#} & \textbf{Axe} & \textbf{Constat et gain} & \textbf{Cout} \\
\midrule
\textbf{R3} & Expirer les sessions en memoire &
Le dictionnaire \code{_sessions} n'a \textbf{ni expiration ni eviction} : toute session creee reste
residente jusqu'au redemarrage. Une expiration apres quelques heures d'inactivite suffit ; l'etat
etant persiste a chaque scene, l'eviction est transparente pour le joueur. Gain : consommation
memoire bornee. & 0,5 j \\
\addlinespace
\textbf{R4} & Verrouiller les versions Python &
Le frontend dispose de \code{package-lock.json} et de \code{npm ci} ; le backend n'a
\textbf{aucun equivalent}. Les bornes hautes du \S\ref{sec:deps} empechent les majeures, mais deux
installations a deux dates peuvent encore differer sur les correctifs. Un
\code{requirements.lock.txt} utilise par le \code{Dockerfile.backend} et la CI supprime cette
derive. & 0,5 j \\
\addlinespace
\textbf{R5} & Limiter le debit sur l'authentification &
La supervision \textbf{detecte} une recherche exhaustive mais rien ne la \textbf{bloque}. Le cout
de bcrypt (219 ms) freine deja l'attaque, mais constitue aussi un \textbf{vecteur de deni de
service} : quelques dizaines de requetes simultanees saturent le processus. Une limitation par IP
traite les deux problemes. & 0,5 j \\
\bottomrule
\end{tabularx}
\normalsize

\subsection{R6 a R8 --- Securisation et hygiene}

\footnotesize
\begin{tabularx}{\linewidth}{@{}L{1.1cm} L{3.4cm} X L{1.3cm}@{}}
\toprule
\textbf{\#} & \textbf{Axe} & \textbf{Constat et gain} & \textbf{Cout} \\
\midrule
\textbf{R6} & Desactiver le mode debug \textcolor{VertOk}{\textbf{(applique)}} &
\code{FastAPI(..., debug=True)} etait ecrit \textbf{en dur} : toute exception non geree exposait la
trace d'execution et des extraits de code source en production (OWASP~A05). \textbf{Applique en
1.3.0} : le mode est pilote par \code{API_DEBUG}, absente par defaut, et un test verifie qu'il
reste desactive. Le cout etant inferieur a l'heure et l'enjeu relevant de la securite, il n'y avait
pas lieu de differer. & < 1 h \\
\addlinespace
\textbf{R7} & Retention sur \code{session_events} &
141 lignes, \textbf{aucun index}, et surtout \textbf{aucune lecture} dans le code : la table est
alimentee par \code{record_event()} mais jamais interrogee. Deux options s'excluent : soit la
donnee est utile et il faut l'exposer et l'indexer sur \code{user_id}, soit elle ne l'est pas et il
faut purger. La conserver sans usage cumule cout de stockage et obligation de protection des
donnees personnelles, sans contrepartie. & 0,5 j \\
\addlinespace
\textbf{R8} & Conteneurs sans privileges root &
Ni \code{Dockerfile.backend} ni \code{frontend/Dockerfile} ne comportent de directive \code{USER}.
Gain : une execution de code arbitraire n'obtient plus root dans le conteneur. Ajuster les droits
sur les volumes \code{/app/data} et \code{/app/saves}. & 0,5 j \\
\bottomrule
\end{tabularx}
\normalsize

\subsection{Priorisation et suivi des gains}

\footnotesize
\begin{tabularx}{\linewidth}{@{}c L{4.6cm} L{2.1cm} X@{}}
\toprule
\textbf{Rang} & \textbf{Recommandation} & \textbf{Delai} & \textbf{Moyen de verification du gain} \\
\midrule
-- & R6 --- desactiver \code{debug=True} & \textcolor{VertOk}{\textbf{Fait (1.3.0)}} & Controle n$^\circ$4 du test de fumee \\
1 & R1 --- chargement 3D a la demande & Immediat & \code{benchmark.py}, section poids livre \\
2 & R5 --- limitation de debit sur le login & Court terme & \code{rpg40k_auth_attempts_total} \\
3 & R3 --- expiration des sessions & Court terme & \code{rpg40k_active_sessions} \\
4 & R2 --- fenetre glissante du contexte & Court terme & \code{rpg40k_gm_context_messages} \\
5 & R4 --- verrouillage des versions Python & Immediat & Comparaison CI / production \\
6 & R7 --- retention \code{session_events} & Moyen terme & \code{benchmark.py}, section base \\
7 & R8 --- conteneurs sans root & Moyen terme & \code{docker compose exec backend id} \\
\bottomrule
\end{tabularx}
\normalsize

\textbf{Charge residuelle : 4 a 5,5 jours-homme}, R6 etant deja appliquee. Aucune refonte
architecturale : toutes les evolutions sont incrementales, deployables par le processus du
\S\ref{sec:cicd}, et \textbf{chacune dispose deja d'un moyen de mesure} (benchmark ou sonde), de
sorte que le gain pourra etre constate plutot que suppose.

\subsection{Contribution a l'attractivite}

\textbf{Rapidite} (R1, R2) : demarrage plus rapide, rythme preserve en partie longue.
\textbf{Fiabilite} (R3, R4) : service stable, comportement identique d'un environnement a l'autre.
\textbf{Confiance} (R5, R6, R8) : compte protege, absence de fuite d'information.
\textbf{Evolutivite} (R7). L'attractivite d'un jeu narratif tient d'abord a la \textbf{fluidite de
l'experience} : un chargement long ou une attente croissante degradent l'engagement bien avant
qu'une fonctionnalite ne manque --- d'ou la priorite de R1 et R2 sur toute extension fonctionnelle.

\couverture
\footnotesize
\begin{tabularx}{\linewidth}{@{}L{4.4cm} X@{}}
\toprule
\textbf{Critere} & \textbf{Traitement} \\
\midrule
Recommandations argumentees, gains evaluables (cout, delai) & Huit axes, chacun avec indicateur mesure, gain chiffre, cout en jours-homme et priorite \\
Realistes et realisables & Aucune refonte ; evolutions incrementales deployables par le \S\ref{sec:cicd} ; reserves de faisabilite (R4, R8) explicitees \\
Renforcent l'attractivite & Quatre axes relies a l'effet percu par le joueur, priorisation justifiee de R1 et R2 \\
\bottomrule
\end{tabularx}
\normalsize

% ===========================================================================
\section{Journal des versions}
\label{sec:journal}

\competence{Competence C4.3.2}{Etablir un journal des versions deployees en y integrant la
documentation des correctifs realises pour suivre les evolutions realisees sur le logiciel.}

\defaut{Defaut identifie --- le logiciel ignorait sa propre version}{%
\code{backend/api.py} declarait \code{version="1.0.0"} en dur alors que le journal en etait a
\code{1.2.0}. Or le formulaire de consignation demande au joueur de relever la valeur retournee par
\code{/api/health} : un signalement aurait donc porte une version \textbf{erronee}, rendant
impossible le rattachement d'une anomalie a un etat du code.}

\textbf{Correctif} : fichier \code{VERSION} a la racine (source unique de verite),
\code{backend/version.py} qui l'expose et lit le journal, et \textbf{7 tests} interdisant la
derive. La protection a ete verifiee en portant volontairement \code{VERSION} a une valeur
incoherente :

\begin{lstlisting}[style=sortie]
$ printf '9.9.9\n' > VERSION && pytest tests/test_version.py -q
FAILED test_la_version_correspond_au_journal
FAILED test_le_journal_documente_la_version_courante        2 failed, 5 passed
\end{lstlisting}

Toute publication oubliant de documenter sa version fait desormais echouer la CI.

\subsection{Exemplaire du journal --- version 1.3.0}

\noindent
\begin{minipage}{\linewidth}
\setlength{\fboxsep}{7pt}\setlength{\fboxrule}{0.8pt}
\fcolorbox{GrisBord}{white}{\begin{minipage}{\dimexpr\linewidth-2\fboxsep-2\fboxrule}
\footnotesize
{\normalsize\bfseries\color{AccentPrim}[1.3.0] --- Maintien en condition operationnelle --- 19/08/2026}\par
\vspace{2pt}{\itshape Ref. branche : \texttt{fix/ANO-2026-001-marqueurs-etat}.}

\vspace{5pt}{\bfseries\color{RougeKo}Corrige}
\begin{itemize}[leftmargin=*,itemsep=2pt,topsep=2pt]
  \item \textbf{\code{ANO-2026-001}} --- perte de la progression d'une scene sur un marqueur d'etat
        non numerique. Le flux SSE s'interrompait avant l'evenement \code{done} \emph{et} avant
        \code{session.save()}. Correctif en defense en profondeur ; 18 tests, 16 en echec avant.
  \item \textbf{\code{ANO-2026-002}} --- perte \emph{partielle} au redemarrage : \code{Session.save()}
        omettait la fiche de personnage. 8 tests, 6 en echec avant.
  \item \textbf{\code{ANO-2026-003}} --- test d'animation intermittent en CI (flottant de haute
        precision). 407 cas sur 300 000. Tests rendus deterministes.
  \item \textbf{Version rapportee incorrecte} et \textbf{documentation du deploiement} contredite
        par le code.
\end{itemize}

\vspace{3pt}{\bfseries\color{OrangeAlerte}Securite}
\begin{itemize}[leftmargin=*,itemsep=1pt,topsep=2pt]
  \item Mode debug desactive par defaut (R6) ; garde-fou \code{JWT_SECRET} au deploiement ; sonde
        et alerte sur les tentatives d'authentification.
\end{itemize}

\vspace{3pt}{\bfseries\color{VertOk}Ajoute}
\begin{itemize}[leftmargin=*,itemsep=1pt,topsep=2pt]
  \item \emph{Supervision} : \code{monitoring.py}, \code{/api/health/ready}, \code{/api/metrics},
        18 regles d'alerte, tableau de bord Grafana.
  \item \emph{Dependances} : \code{dependabot.yml} (4 ecosystemes), bornes hautes,
        \code{check_updates.py}.
  \item \emph{Deploiement} : \code{smoke_test.sh} (5 controles), retour arriere automatique.
  \item \emph{Qualite} : formulaires de consignation, \code{benchmark.py}, fichier \code{VERSION}.
\end{itemize}

\vspace{3pt}{\bfseries\color{AccentSec}Tests} --- \code{test_monitoring.py} (20),
\code{test_state_markers.py} (18), \code{test_version.py} (7),
\code{test_character_persistence.py} (8). \textbf{Total : 82 $\rightarrow$ 135 au vert}, plus 30
tests frontend.
\end{minipage}}
\end{minipage}

\vspace{5pt}
\textbf{Regles de tenue} --- une section par version, la plus recente en premier ; rubriques
normalisees (\emph{Ajoute}, \emph{Modifie}, \emph{Corrige}, \emph{Securite}, \emph{Tests}) ; chaque
correctif nomme l'anomalie, son symptome, sa cause et son correctif ; chaque version reference son
tag, ce qui permet le retour arriere ; les rubriques \emph{Securite} sont toujours explicites. Une
entree est redigee \textbf{au moment du commit}, jamais reconstituee a la publication : c'est la
seule facon d'obtenir une description fidele.

La chaine est verifiable \textbf{dans les deux sens} : depuis une version affichee par
\code{/api/health} on retrouve l'entree du journal, le commit puis la fiche d'anomalie ;
inversement, une fiche mene au commit qui la corrige et a la version qui l'a deployee.

\couverture
\footnotesize
\begin{tabularx}{\linewidth}{@{}L{4.4cm} X@{}}
\toprule
\textbf{Critere} & \textbf{Traitement} \\
\midrule
Le journal contient les ameliorations de la version & Exemplaire 1.3.0 : \emph{Corrige} (4), \emph{Securite} (3), \emph{Ajoute} (4 domaines), \emph{Tests} \\
Les correctifs deployes sont documentes & Chaque correctif nomme anomalie, symptome, cause racine, correctif et couverture de tests ; correspondance versions / preuves \\
\bottomrule
\end{tabularx}
\normalsize

% ===========================================================================
\section{Probleme resolu avec le support client}
\label{sec:support}

\competence{Competence C4.3.3}{Collaborer avec les equipes de support, en fournissant une expertise
technique, en repondant aux retours clients, en resolvant des problemes complexes afin d'ameliorer
le logiciel.}

\reserve{\textbf{Precision de contexte.} Le projet est developpe par une personne seule : il
n'existe pas d'equipe de support distincte. Les roles decrits correspondent aux \textbf{fonctions
reellement exercees} --- reception et qualification d'une part, expertise technique d'autre part ---
et non a des personnes distinctes. L'anomalie, son diagnostic et son correctif sont authentiques.
Inventer des echanges entre equipes fictives aurait retire toute valeur probante a cet exemple.}

\subsection{Contexte du retour et probleme a resoudre}

\textbf{Signalement} : « Mon personnage est revenu a zero mais j'ai garde mon niveau et mon
equipement. » Gravite majeure, perimetre sauvegarde, version 1.2.0, frequence « deux fois, mais pas
a chaque fois ». Le joueur est niveau 3 avec tout son equipement, mais son personnage est de
nouveau indemne, ses rations pleines, ses points d'attribut disparus.

Trois caracteristiques en font un cas difficile. \textbf{1.} La perte est \textbf{partielle}, donc
contre-intuitive : une perte totale orienterait vers l'ecriture disque, tandis que cette
incoherence apparente suggere a tort un probleme d'affichage. \textbf{2.} Le symptome parait
\textbf{intermittent} : la formulation « pas a chaque fois » est exacte mais decrit une correlation
que le joueur ne peut pas percevoir. \textbf{3.} Le contournement evident est \textbf{inoperant} :
sauvegarder manuellement ne change rien.

\subsection{Contribution des parties prenantes}

\textbf{Niveau 1 --- reception.} Trois apports decisifs obtenus par le formulaire structure : la
\textbf{distinction precise} entre donnees perdues et conservees, qui orientera tout le diagnostic ;
la \textbf{qualification en gravite majeure}, qui declenche une prise en charge sous 48 h meme sans
blocage ; et surtout la \textbf{reformulation de l'intermittence} --- interroge sur ce qui
distinguait les sessions touchees, le joueur precise qu'il revenait « le lendemain », par
opposition a un simple rechargement. Cette precision transforme un symptome « aleatoire » en
condition reproductible : un delai suffisant pour que le serveur ait redemarre. Sans elle, le
niveau 2 aurait cherche une anomalie d'ecriture disque, alors que le defaut se situe au
\textbf{rechargement}.

\textbf{Niveau 2 --- expertise.} L'hypothese est reproduite en vidant les sessions residentes, ce
qui simule un redemarrage. Le releve reproduit \emph{exactement} la repartition decrite :

\begin{lstlisting}[style=sortie]
donnee                     avant redemarrage    apres     verdict
rations (personnage)                       1        3     >>> PERDU <<<
blessures (personnage)                 Grave  Indemne     >>> PERDU <<<
points d'attribut                          9   remis a 0  >>> PERDU <<<
niveau / xp (progression)            3 / 250    3 / 250    CONSERVE
\end{lstlisting}

\defaut{Cause racine --- un sous-systeme absent de la liste}{%
\code{Session.save()} persiste \emph{six} sous-systemes (progression, inventaire, carte, quetes,
relations, equipe) mais omet \code{self.character}. Au demarrage, la fiche est donc reconstruite
depuis le modele vierge. \textbf{La liste des donnees perdues correspond exactement au contenu de
\code{CharacterState}, celle des donnees conservees aux six sous-systemes persistes : la
repartition observee par le joueur est la signature directe de la cause.}}

L'examen confirme qu'il s'agit d'un oubli et non d'un choix : la classe possede \code{to_dict()} et
\code{save()}, mais \textbf{aucune reciproque \code{from_dict()}}. L'etat pouvait etre serialise mais
jamais relu --- l'omission n'avait donc jamais pu etre detectee, faute de chemin de lecture. Aucun
test ne simulait de redemarrage : le defaut n'etait observable qu'au rechargement a froid, soit
exactement la condition que le joueur avait decrite sans pouvoir la nommer.

\subsection{Resolution apportee}

Trois volets : ajout de la reciproque \code{CharacterState.from_dict()} ; persistance de
\code{character.yaml} dans \code{Session.save()} ; rechargement avec \textbf{repli} sur le modele
vierge pour les parties anterieures au correctif et en cas de fichier illisible. Cette
compatibilite ascendante est une exigence du support : un correctif rendant inutilisables les
parties existantes transformerait une perte partielle en \textbf{perte totale}.

\textbf{8 tests de non-regression} : \textcolor{RougeKo}{\textbf{6 en echec}} avant correctif,
\textcolor{VertOk}{\textbf{8 au vert}} apres. Les deux qui passaient deja sont ceux de
compatibilite ascendante, que le correctif ne devait pas modifier. Un test verifie specifiquement
la \textbf{coherence} entre fiche et progression, puisque c'est leur divergence qui constituait le
symptome.

\textbf{Reponse au joueur} : diagnostic en termes non techniques, explication de l'intermittence,
contournement provisoire, version de livraison. Le deuxieme point est le plus important : il
\textbf{valide l'observation du joueur} au lieu de la contredire. Un signalement qualifie
d'« irreproductible » decourage les suivants, qui sont pourtant la principale source de detection
sur un logiciel a faible base d'utilisateurs.

\subsection{Amelioration au-dela du correctif}

Une \textbf{couverture de test manquante} a ete comblee (aucun test ne simulait de redemarrage :
c'est desormais le cas pour \emph{tous} les sous-systemes persistes) ; une \textbf{asymetrie d'API}
a ete corrigee (\code{to_dict()} sans \code{from_dict()}), cause profonde de l'omission ; et le
diagnostic a confirme l'interet de la recommandation R3.

\couverture
\footnotesize
\begin{tabularx}{\linewidth}{@{}L{4.4cm} X@{}}
\toprule
\textbf{Critere} & \textbf{Traitement} \\
\midrule
Contexte du retour et probleme a resoudre & Fiche telle que recue, puis analyse : perte partielle contre-intuitive, intermittence apparente, contournement inoperant \\
La resolution apportee & Trois volets, compatibilite ascendante, verification par 8 tests executes avant et apres, reponse formulee au joueur \\
Contribution des parties prenantes & Apport detaille de chaque fonction, en indiquant ce qui aurait manque sans elle \\
\bottomrule
\end{tabularx}
\normalsize

% ===========================================================================
\section{Conclusion}

Le dispositif presente ici a ete \textbf{construit}, non decrit a partir de l'existant. Le projet
disposait au depart d'une chaine d'integration continue et d'un journal des versions, mais d'aucune
supervision, d'aucun outillage de veille, d'aucun processus de collecte d'anomalies et d'aucune
validation post-deploiement.

\footnotesize
\begin{tabularx}{\linewidth}{@{}L{4.2cm} L{4.6cm} X@{}}
\toprule
\textbf{Domaine} & \textbf{Etat initial} & \textbf{Etat final} \\
\midrule
Gestion des dependances & Planchers seuls, majeures subies & Bornes hautes, Dependabot hebdomadaire, veille en CI \\
Supervision & Aucune ; \code{/api/health} n'interrogeait rien & 19 sondes, sonde d'aptitude reelle, 18 regles, tableau de bord \\
Collecte des anomalies & Aucun outil & 3 formulaires structures, saisie libre desactivee \\
Validation du deploiement & \code{curl} sur la sonde de vivacite & 5 controles fonctionnels, rollback automatique \\
Indicateurs de performance & Aucune mesure & \code{benchmark.py} reproductible \\
Version rapportee & \code{1.0.0} en dur, journal a \code{1.2.0} & Source unique de verite, coherence testee \\
Tests backend & 82 & \textbf{135} \\
\bottomrule
\end{tabularx}
\normalsize

\vspace{5pt}
\textbf{Trois anomalies reelles} ont ete decouvertes, reproduites, corrigees et couvertes par des
tests qui echouent sur le code anterieur --- condition sans laquelle un test de non-regression ne
prouve rien. Les deux premieres partagent un trait instructif : dans les deux cas, \textbf{une
protection existait deja ailleurs dans le meme fichier}. La branche \code{tracks} interceptait
l'exception que la branche \code{resources} laissait passer ; six sous-systemes etaient persistes
quand le septieme etait omis. Le defaut ne venait pas d'une meconnaissance du probleme, mais d'une
\textbf{application incomplete d'une solution deja connue} --- ce qui plaide pour des tests portant
sur la \emph{symetrie} des traitements, et non seulement sur les chemins nominaux.

La charge residuelle des recommandations non appliquees est estimee a \textbf{4 a 5,5 jours-homme},
sans refonte architecturale. Les deux axes prioritaires portent tous deux sur la \textbf{fluidite
percue}, qui conditionne l'engagement davantage que l'ajout de fonctionnalites. Le dispositif rend
ces evolutions \textbf{mesurables} : chaque recommandation dispose d'un moyen de verification deja
operationnel, de sorte que le gain pourra etre constate plutot que suppose.

% ===========================================================================
\section{Annexe A --- Artefacts ajoutes au depot}

Chaque element du dossier correspond a un fichier verifiable. Le tableau ci-dessous recense les
artefacts crees ou modifies, avec la competence qu'ils servent.

\footnotesize
\begin{tabularx}{\linewidth}{@{}L{5.6cm} c X@{}}
\toprule
\textbf{Fichier} & \textbf{Etat} & \textbf{Role} \\
\midrule
\multicolumn{3}{@{}l}{\textbf{\color{AccentSec}Gestion des dependances --- C4.1.1}} \\
\code{.github/dependabot.yml} & cree & Surveillance hebdomadaire des 4 ecosystemes \\
\code{requirements.txt} & modifie & Bornes hautes interdisant les montees majeures \\
\code{scripts/check_updates.py} & cree & Rapport de veille classant les obsolescences \\
\code{.github/workflows/ci.yml} & modifie & Veille integree au job \code{security-audit} \\
\midrule
\multicolumn{3}{@{}l}{\textbf{\color{AccentSec}Supervision --- C4.1.2}} \\
\code{backend/monitoring.py} & cree & Sondes, middleware, journalisation, sonde d'aptitude \\
\code{backend/api.py} & modifie & \code{/api/health/ready}, \code{/api/metrics}, instrumentation \\
\code{monitoring/prometheus.yml} & cree & Collecte, cibles, auto-supervision \\
\code{monitoring/alert_rules.yml} & cree & 18 regles reparties en 6 groupes \\
\code{monitoring/alertmanager.yml} & cree & Routage par severite, groupement, inhibition \\
\code{monitoring/grafana/} & cree & Tableau de bord provisionne (13 panneaux) \\
\code{docker-compose.monitoring.yml} & cree & Pile de supervision en surcouche \\
\midrule
\multicolumn{3}{@{}l}{\textbf{\color{AccentSec}Anomalies et support --- C4.2.1, C4.3.3}} \\
\code{.github/ISSUE_TEMPLATE/} & cree & 3 formulaires, saisie libre desactivee \\
\code{src/state.py} & modifie & Correctifs \code{ANO-2026-001} et \code{ANO-2026-002} \\
\midrule
\multicolumn{3}{@{}l}{\textbf{\color{AccentSec}Deploiement --- C4.2.2}} \\
\code{scripts/smoke_test.sh} & cree & 5 controles post-deploiement \\
\code{.github/workflows/deploy-vps.yml} & modifie & Test de fumee, rollback, garde-fou \code{JWT_SECRET} \\
\midrule
\multicolumn{3}{@{}l}{\textbf{\color{AccentSec}Performance et versions --- C4.3.1, C4.3.2}} \\
\code{scripts/benchmark.py} & cree & Mesure des indicateurs de performance \\
\code{VERSION} & cree & Source unique de verite du numero de version \\
\code{backend/version.py} & cree & Lecture de la version et du journal \\
\code{CHANGELOG.md} & modifie & Entree \code{[1.3.0]} documentant les correctifs \\
\code{frontend/scripts/capture_screenshots.mjs} & cree & Captures d'ecran reproductibles de ce dossier \\
\bottomrule
\end{tabularx}
\normalsize

% ===========================================================================
\section{Annexe B --- Inventaire des tests}

La suite backend passe de \textbf{82 a 135 tests}. Les quatre fichiers ajoutes couvrent
exactement les dispositifs et les anomalies decrits dans ce dossier.

\footnotesize
\begin{tabularx}{\linewidth}{@{}L{5.0cm} c c X@{}}
\toprule
\textbf{Fichier de test} & \textbf{Tests} & \textbf{Avant correctif} & \textbf{Ce qu'il garantit} \\
\midrule
\code{tests/test_monitoring.py} & 20 & --- &
Normalisation des libelles (cardinalite bornee), comptage du trafic par famille de code, sondes de
securite, et surtout : la sonde d'aptitude \textbf{detecte reellement} une base corrompue et un
fichier de jeu absent. \\
\addlinespace
\code{tests/test_state_markers.py} & 18 & \textcolor{RougeKo}{\textbf{16 en echec}} &
Aucune valeur textuelle ne leve d'exception ; un marqueur rejete ne corrompt pas l'etat ; un
marqueur invalide n'annule pas ses voisins valides ; le flux SSE emet \code{done} et
\code{session.save()} est atteint. \\
\addlinespace
\code{tests/test_version.py} & 7 & --- &
Coherence entre \code{VERSION}, la valeur exposee par \code{/api/health} et l'entree la plus
recente du journal ; mode debug desactive par defaut. \\
\addlinespace
\code{tests/test_character_persistence.py} & 8 & \textcolor{RougeKo}{\textbf{6 en echec}} &
Ressources, blessures et points d'attribut survivent a un redemarrage ; coherence fiche /
progression ; compatibilite ascendante avec les sauvegardes anterieures et repli sur fichier
illisible. \\
\midrule
\textbf{Total ajoute} & \textbf{53} & \textbf{22 en echec} & \\
\bottomrule
\end{tabularx}
\normalsize

\vspace{5pt}
La colonne \emph{Avant correctif} est la plus importante : elle atteste que les tests de
non-regression \textbf{echouent effectivement} sur le code anterieur. Un test qui passerait dans les
deux etats ne prouverait rien de l'anomalie qu'il pretend couvrir.

\vspace{5pt}
Cote frontend, les \textbf{30 tests} existants sont conserves ; cinq d'entre eux ont ete rendus
\textbf{deterministes} au titre de \code{ANO-2026-003}, sans modification du code de production.
